\documentclass[runningheads]{llncs}

\usepackage{booktabs}
\usepackage{listings}
\usepackage{url}

\input{coq-listings}

\lstset{
  basicstyle=\ttfamily\small,
  columns=fullflexible,
  keepspaces=true,
  frame=single,
  xleftmargin=0.5em,
  xrightmargin=0.5em
}

\IfFileExists{generated/results.tex}
  {% Generated by scripts/make_paper_results.py.
% Re-run the artifact workflow before submitting final numbers.
\newcommand{\JsonifiableDiscovered}{342}
\newcommand{\JsonifiableProbed}{311}
\newcommand{\JsonifiableSkipped}{31}
\newcommand{\JsonifiableSkippedParameterizedModule}{24}
\newcommand{\JsonifiableSkippedModuleType}{7}
\newcommand{\JsonifiableProbeSuccesses}{81}
\newcommand{\JsonifiableProbeFailures}{230}
\newcommand{\JsonifiableProbeTimeouts}{0}
\newcommand{\JsonifiableTimedDerivations}{81}
\newcommand{\JsonifiableBenchmarkWallSeconds}{5.21}
\newcommand{\JsonifiableDerivationSumSeconds}{4.29}
\newcommand{\JsonifiableDerivationMeanSeconds}{0.053}
\newcommand{\JsonifiableDerivationMedianSeconds}{0.034}
\newcommand{\JsonifiableDerivationPNineZeroSeconds}{0.112}
\newcommand{\JsonifiableDerivationMaxSeconds}{0.521}
\newcommand{\JsonifiableSuccessInductives}{50}
\newcommand{\JsonifiableSuccessVariants}{17}
\newcommand{\JsonifiableSuccessRecords}{14}
\newcommand{\JsonifiableSuccessCorelib}{39}
\newcommand{\JsonifiableSuccessStdlib}{42}
\newcommand{\JsonifiableExtractionGeneratedRuns}{10}
\newcommand{\JsonifiableExtractionGeneratedMean}{0.055843}
\newcommand{\JsonifiableExtractionGeneratedMedian}{0.055783}
\newcommand{\JsonifiableExtractionHandwrittenRuns}{10}
\newcommand{\JsonifiableExtractionHandwrittenMean}{0.055628}
\newcommand{\JsonifiableExtractionHandwrittenMedian}{0.055635}

\newcommand{\JsonifiableFailureRows}{%
prop-target-not-supported & 150 \\
dependent-field-not-supported & 48 \\
indexed-family-not-supported & 21 \\
function-field-not-supported & 4 \\
proof-field-not-supported & 3 \\
sort-field-not-supported & 3 \\
coinductive-not-supported & 1 \\
}

\newcommand{\JsonifiableSkippedRows}{%
parameterized-module-scope & 24 & The declaration appears inside a parameterized module/functor body, so there is no closed global reference to derive until the functor is instantiated. \\
module-type-scope & 7 & The declaration appears in a module type/signature, which describes an interface component rather than a closed inductive declaration. \\
}

\newcommand{\JsonifiableSlowestRows}{%
Corelib.Init.Byte.byte & Inductive & 0.521 \\
Stdlib.rtauto.Rtauto.proof & Inductive & 0.182 \\
Stdlib.btauto.Reflect.formula & Inductive & 0.174 \\
Stdlib.micromega.ZMicromega.ZArithProof & Inductive & 0.166 \\
Corelib.Floats.FloatClass.float\_class & Variant & 0.142 \\
Stdlib.micromega.RingMicromega.Psatz & Inductive & 0.137 \\
Stdlib.Sorting.Heap.Tree & Inductive & 0.125 \\
Stdlib.setoid\_ring.Field\_theory.FExpr & Inductive & 0.116 \\
}

\newcommand{\JsonifiableFailureExampleRows}{%
prop-target-not-supported & 150 & Corelib.Classes.CMorphisms.apply\_subrelation & Targets live in Prop, while Jsonifiable builds computational data in Type. \\
dependent-field-not-supported & 48 & Corelib.Init.Specif.sig & A constructor field type depends on an earlier constructor argument, so a plain JSON decoder cannot reconstruct the dependent evidence. \\
indexed-family-not-supported & 21 & Corelib.Init.Datatypes.reflect & The declaration is a family indexed by values, not a plain data type after parameters. \\
function-field-not-supported & 4 & Corelib.ssr.ssrbool.implies & A constructor field is function-valued, and there is no generic JSON encoding for arbitrary functions. \\
proof-field-not-supported & 3 & Corelib.Init.Specif.sumbool & A constructor field lives in Prop or SProp, so the ordinary data round-trip contract would require proof erasure or proof reconstruction. \\
sort-field-not-supported & 3 & Corelib.Classes.RelationClasses.Tlist & A constructor field is itself a sort such as Type or Set, while Jsonifiable serializes data values rather than type-level objects. \\
coinductive-not-supported & 1 & Stdlib.Streams.Streams.Stream & The declaration is potentially infinite data, while the deriver targets inductive data. \\
}

\newcommand{\JsonifiableExtractionRows}{%
enum256\_generated & 10 & 0.055843 & 0.055783 & 0.052521 & 0.059224 \\
enum256\_handwritten & 10 & 0.055628 & 0.055635 & 0.053497 & 0.059411 \\
\midrule
point2d\_generated & 10 & 0.001283 & 0.001263 & 0.001122 & 0.001663 \\
point2d\_handwritten & 10 & 0.001378 & 0.001371 & 0.001300 & 0.001454 \\
\midrule
userrole\_generated & 10 & 0.002316 & 0.002353 & 0.002034 & 0.002450 \\
userrole\_handwritten & 10 & 0.003285 & 0.003339 & 0.003051 & 0.003475 \\
\midrule
serverconfig\_generated & 10 & 0.000749 & 0.000754 & 0.000708 & 0.000802 \\
serverconfig\_handwritten & 10 & 0.001961 & 0.001922 & 0.001869 & 0.002109 \\
\midrule
instruction\_generated & 10 & 0.002213 & 0.002225 & 0.002010 & 0.002359 \\
instruction\_handwritten & 10 & 0.003491 & 0.003402 & 0.003229 & 0.004540 \\
\midrule
bintree\_generated & 10 & 0.004180 & 0.004129 & 0.004043 & 0.004450 \\
bintree\_handwritten & 10 & 0.006416 & 0.006408 & 0.006149 & 0.006721 \\
}
}
  {% Generated by scripts/make_paper_results.py.
% Re-run the artifact workflow before submitting final numbers.
\newcommand{\JsonifiableDiscovered}{342}
\newcommand{\JsonifiableProbed}{311}
\newcommand{\JsonifiableSkipped}{31}
\newcommand{\JsonifiableSkippedParameterizedModule}{24}
\newcommand{\JsonifiableSkippedModuleType}{7}
\newcommand{\JsonifiableProbeSuccesses}{81}
\newcommand{\JsonifiableProbeFailures}{230}
\newcommand{\JsonifiableProbeTimeouts}{0}
\newcommand{\JsonifiableTimedDerivations}{81}
\newcommand{\JsonifiableBenchmarkWallSeconds}{5.21}
\newcommand{\JsonifiableDerivationSumSeconds}{4.29}
\newcommand{\JsonifiableDerivationMeanSeconds}{0.053}
\newcommand{\JsonifiableDerivationMedianSeconds}{0.034}
\newcommand{\JsonifiableDerivationPNineZeroSeconds}{0.112}
\newcommand{\JsonifiableDerivationMaxSeconds}{0.521}
\newcommand{\JsonifiableSuccessInductives}{50}
\newcommand{\JsonifiableSuccessVariants}{17}
\newcommand{\JsonifiableSuccessRecords}{14}
\newcommand{\JsonifiableSuccessCorelib}{39}
\newcommand{\JsonifiableSuccessStdlib}{42}
\newcommand{\JsonifiableExtractionGeneratedRuns}{10}
\newcommand{\JsonifiableExtractionGeneratedMean}{0.055843}
\newcommand{\JsonifiableExtractionGeneratedMedian}{0.055783}
\newcommand{\JsonifiableExtractionHandwrittenRuns}{10}
\newcommand{\JsonifiableExtractionHandwrittenMean}{0.055628}
\newcommand{\JsonifiableExtractionHandwrittenMedian}{0.055635}

\newcommand{\JsonifiableFailureRows}{%
prop-target-not-supported & 150 \\
dependent-field-not-supported & 48 \\
indexed-family-not-supported & 21 \\
function-field-not-supported & 4 \\
proof-field-not-supported & 3 \\
sort-field-not-supported & 3 \\
coinductive-not-supported & 1 \\
}

\newcommand{\JsonifiableSkippedRows}{%
parameterized-module-scope & 24 & The declaration appears inside a parameterized module/functor body, so there is no closed global reference to derive until the functor is instantiated. \\
module-type-scope & 7 & The declaration appears in a module type/signature, which describes an interface component rather than a closed inductive declaration. \\
}

\newcommand{\JsonifiableSlowestRows}{%
Corelib.Init.Byte.byte & Inductive & 0.521 \\
Stdlib.rtauto.Rtauto.proof & Inductive & 0.182 \\
Stdlib.btauto.Reflect.formula & Inductive & 0.174 \\
Stdlib.micromega.ZMicromega.ZArithProof & Inductive & 0.166 \\
Corelib.Floats.FloatClass.float\_class & Variant & 0.142 \\
Stdlib.micromega.RingMicromega.Psatz & Inductive & 0.137 \\
Stdlib.Sorting.Heap.Tree & Inductive & 0.125 \\
Stdlib.setoid\_ring.Field\_theory.FExpr & Inductive & 0.116 \\
}

\newcommand{\JsonifiableFailureExampleRows}{%
prop-target-not-supported & 150 & Corelib.Classes.CMorphisms.apply\_subrelation & Targets live in Prop, while Jsonifiable builds computational data in Type. \\
dependent-field-not-supported & 48 & Corelib.Init.Specif.sig & A constructor field type depends on an earlier constructor argument, so a plain JSON decoder cannot reconstruct the dependent evidence. \\
indexed-family-not-supported & 21 & Corelib.Init.Datatypes.reflect & The declaration is a family indexed by values, not a plain data type after parameters. \\
function-field-not-supported & 4 & Corelib.ssr.ssrbool.implies & A constructor field is function-valued, and there is no generic JSON encoding for arbitrary functions. \\
proof-field-not-supported & 3 & Corelib.Init.Specif.sumbool & A constructor field lives in Prop or SProp, so the ordinary data round-trip contract would require proof erasure or proof reconstruction. \\
sort-field-not-supported & 3 & Corelib.Classes.RelationClasses.Tlist & A constructor field is itself a sort such as Type or Set, while Jsonifiable serializes data values rather than type-level objects. \\
coinductive-not-supported & 1 & Stdlib.Streams.Streams.Stream & The declaration is potentially infinite data, while the deriver targets inductive data. \\
}

\newcommand{\JsonifiableExtractionRows}{%
enum256\_generated & 10 & 0.055843 & 0.055783 & 0.052521 & 0.059224 \\
enum256\_handwritten & 10 & 0.055628 & 0.055635 & 0.053497 & 0.059411 \\
\midrule
point2d\_generated & 10 & 0.001283 & 0.001263 & 0.001122 & 0.001663 \\
point2d\_handwritten & 10 & 0.001378 & 0.001371 & 0.001300 & 0.001454 \\
\midrule
userrole\_generated & 10 & 0.002316 & 0.002353 & 0.002034 & 0.002450 \\
userrole\_handwritten & 10 & 0.003285 & 0.003339 & 0.003051 & 0.003475 \\
\midrule
serverconfig\_generated & 10 & 0.000749 & 0.000754 & 0.000708 & 0.000802 \\
serverconfig\_handwritten & 10 & 0.001961 & 0.001922 & 0.001869 & 0.002109 \\
\midrule
instruction\_generated & 10 & 0.002213 & 0.002225 & 0.002010 & 0.002359 \\
instruction\_handwritten & 10 & 0.003491 & 0.003402 & 0.003229 & 0.004540 \\
\midrule
bintree\_generated & 10 & 0.004180 & 0.004129 & 0.004043 & 0.004450 \\
bintree\_handwritten & 10 & 0.006416 & 0.006408 & 0.006149 & 0.006721 \\
}
}

\title{Engineering Certified JSON Derivers for Rocq with ELPI}
\titlerunning{Certified JSON Derivers for Rocq}

\author{Will Thomas}
\authorrunning{W. Thomas}
\institute{University of Kansas}

\begin{document}
\maketitle

\begin{abstract}
  Rocq developments often need executable serializers for data exchange, testing,
  and extracted applications, but hand-written serializers are tedious and easy to
  disconnect from their intended inverse.  We present an ELPI-based deriver for
  \texttt{rocq-json} that automatically generates \texttt{Jsonifiable} instances
  for a practical fragment of Rocq inductive and record types.  Each generated
  instance contains an encoder, a decoder, and a proof that decoding the encoded
  value succeeds with the original value.  The implementation is not intended as a
  new theory of generic programming, nor does the underlying JSON library aim to
  be an RFC-complete JSON implementation.  Instead, the contribution is a
  performance-conscious and extraction-friendly engineering pattern for mapping
  Rocq data to a JSON-like interchange format while preserving an explicit
  invertibility theorem.  The accompanying artifact rebuilds the tool, probes
Corelib and Stdlib declarations, classifies unsupported cases, compares
generated and handwritten Rocq definitions after extraction, and
regenerates the result tables used by this paper.
\end{abstract}

\section{Introduction}

Serialization code sits at an awkward boundary in proof assistant projects.  It
is usually mundane enough that users do not want to write it by hand, but it is
important enough that a bug can invalidate testing infrastructure or extracted
tools.  In Rocq, this tension is especially visible for data types that are
defined primarily for proofs and then reused in extracted programs.

\texttt{rocq-json} addresses this problem with the following typeclass:

\begin{lstlisting}[language=Coq]
Class Jsonifiable (A : Type) := {
  to_JSON     : A -> JSON;
  from_JSON   : JSON -> Result A string;
  canonical_jsonification : forall (a : A), from_JSON (to_JSON a) = res a
}.
\end{lstlisting}

The deriver described here generates all three components for many ordinary
inductive and record types.  It is implemented in coq-elpi and exposed through
the command \texttt{Elpi derive.jsonifiable T}.

This is a tool paper about useful engineering rather than a claim of novel
foundations.  The interesting part is the integration: ELPI inspection of Rocq
declarations, generated code that remains suitable for extraction, typeclass
parameters for nested serializable fields, and automatic proof construction for
the round-trip theorem.  The same ingredients are broadly useful for other
derivers that construct computational content and accompanying correctness
proofs.

% Reference ideas, intentionally not instantiated here:
% - coq-elpi / ELPI papers and documentation.
% - Rocq/Coq reference or system papers.
% - Prior derive/generic-programming tools in Coq/Rocq.
% - Extraction references.
% - Tool paper examples from verification conferences.

\section{Tool Overview}

The user-facing interface is deliberately small.  Given a declaration such as an
enumeration, algebraic data type, parameterized container, or record, the user
invokes the deriver and receives a registered typeclass instance.  Nullary
constructors are encoded as JSON strings.  Constructors with arguments are
encoded as singleton JSON objects whose key is the constructor name and whose
value stores named fields.  Records are encoded as JSON objects keyed by
projection names.

For example, the following declaration and command are enough to produce an
encoder, decoder, proof, and registered instance:

\begin{lstlisting}[language=Coq]
Inductive UserRole : Type :=
| Admin
| Moderator (level : nat)
| StandardUser (id : nat) (username : string)
| Guest.
Elpi derive.jsonifiable UserRole.
\end{lstlisting}

The generated representation uses a string for nullary constructors and an
object for constructors with fields:

\begin{lstlisting}[language=Coq]
Example role_roundtrip (r : UserRole) :
  from_JSON (to_JSON r) = res r := canonical_jsonification r.

Example moderator_json :
  to_JSON (Moderator 2) =
    JSON_Object [("Moderator", JSON_Object [("level", JSON_Nat 2)])] :=
  eq_refl.
\end{lstlisting}

For parameterized data types, the generated instance abstracts over
\texttt{Jsonifiable} instances for the type parameters.  For recursive types,
the generated encoder and decoder use structurally recursive definitions.  The
current implementation also supports a conservative collection of nested
recursive occurrences through standard data containers such as lists, options,
and products.  More general nested or mutually recursive types are reported as
unsupported rather than handled partially by ad hoc guesses.

\begin{lstlisting}[language=Coq]
Inductive BinTree (A : Type) : Type :=
| BinLeaf
| BinNode (left : BinTree A) (value : A) (right : BinTree A).
Elpi derive.jsonifiable BinTree.
Check BinTree_Jsonifiable :
  forall A, Jsonifiable A -> Jsonifiable (BinTree A).
\end{lstlisting}

The generated proof is part of the public contract.  A successful derivation
does not merely produce functions with plausible behavior; it registers a
\texttt{canonical\_jsonification} proof.  This property is intentionally
one-sided.  It says that values produced by the encoder can be decoded back to
the original value.  It does not claim that every accepted JSON value is in a
canonical printed form.

\section{Implementation}

The deriver follows the shape of the Rocq declaration.  ELPI first reads the
inductive or record view, classifies parameters and constructor arguments, and
then builds separate skeletons for the encoder, decoder, proof, and final
typeclass instance.  The skeletons are elaborated by Rocq before being added as
constants.  This staging has been useful in practice: ELPI performs structural
analysis and code generation, while Rocq remains responsible for checking the
terms that become part of the trusted development.

Two engineering choices are central.  First, field serializers are kept
abstract through typeclass arguments where possible.  This makes generated
instances compose with existing hand-written or derived instances, including
standard containers.  Second, decoder generation is performance-sensitive.
Large enumerations are decoded using a string decision tree that extracts to
direct pattern matching over characters, avoiding a long chain of string
equality tests.  The resulting extracted code is close to the shape a careful
user would write by hand.

The proof-generation strategy is mixed.  General recursive and record cases use
Rocq proof automation specialized to the generated encoders and decoders.  Pure
enumerations use a direct generated proof term, avoiding proof-search overhead
on a common stress case.  This division is not theoretically deep, but it
matters for predictable tool latency.

\section{Evaluation}

We evaluate two questions: how broadly the deriver applies to existing Rocq
declarations, and whether generated extracted code is competitive with
hand-written code on a stress case.

\subsection{Applicability}

The artifact includes a benchmark script that scans installed Corelib and
Stdlib sources for inductive-like declarations, probes each candidate in
isolation, records diagnostics, and then compiles a single benchmark file
containing all successful derivations.  This workflow is deliberately
conservative: unsupported declarations remain visible as categorized failures.

\begin{table}
  \centering
  \caption{Corelib/Stdlib applicability summary.}
  \begin{tabular}{lr}
    \toprule
    Metric                       & Count                        \\
    \midrule
    Discovered declarations      & \JsonifiableDiscovered       \\
    Probed declarations          & \JsonifiableProbed           \\
    Successful probes            & \JsonifiableProbeSuccesses   \\
    Probe failures               & \JsonifiableProbeFailures    \\
    Probe timeouts               & \JsonifiableProbeTimeouts    \\
    Timed successful derivations & \JsonifiableTimedDerivations \\
    \bottomrule
  \end{tabular}
\end{table}

\begin{table}
  \centering
  \caption{Successful derivations by declaration shape and library.}
  \begin{tabular}{lr@{\qquad}lr}
    \toprule
    Shape            & Count                         & Library & Count                      \\
    \midrule
    Inductive        & \JsonifiableSuccessInductives & Corelib & \JsonifiableSuccessCorelib \\
    Variant          & \JsonifiableSuccessVariants   & Stdlib  & \JsonifiableSuccessStdlib  \\
    Record/Structure & \JsonifiableSuccessRecords    &         &                            \\
    \bottomrule
  \end{tabular}
\end{table}

\begin{table}
  \centering
  \caption{Failure categories reported by the probe benchmark.}
  \begin{tabular}{lr}
    \toprule
    Category & Count \\
    \midrule
    \JsonifiableFailureRows
    \bottomrule
  \end{tabular}
\end{table}

The category names in this table are not meant to declare every failure
``invalid input''.  They are a triage device.  Some categories are semantic
non-targets for this version of the tool, such as propositions or coinductive
streams.  Others are engineering limitations or possible tool failures that
should be treated as future work.  Table~\ref{tab:failure-examples} gives one
representative declaration for each category and the rationale used by the
artifact.

\begin{table}
  \centering
  \caption{Representative failure examples and classification rationale.}
  \label{tab:failure-examples}
  \begin{tabular}{lrp{0.27\linewidth}p{0.36\linewidth}}
    \toprule
    Category & Count & Example & Rationale \\
    \midrule
    \JsonifiableFailureExampleRows
    \bottomrule
  \end{tabular}
\end{table}

The important distinction is between unsupported-by-design and unresolved.  The
paper claims the former only for cases where the current \texttt{Jsonifiable}
interface has no clear computational target, for example \texttt{Prop} targets
or potentially infinite coinductive data.  Categories such as
\texttt{to-json-elaboration-failed}, \texttt{to-json-generation-not-supported},
and \texttt{scanner-or-logical-path} are better read as limitations of the
implementation or benchmark harness.  They are reported explicitly so the
success rate is auditable rather than inflated by silent filtering.

\subsection{Derivation Time}

The combined successful benchmark records Rocq's \texttt{Time} output for each
derived instance.  In the generated run used for this draft, the combined
benchmark compiled successfully with wall-clock time
\JsonifiableBenchmarkWallSeconds{} seconds; the sum of Rocq-reported derivation
transactions was \JsonifiableDerivationSumSeconds{} seconds.  The mean
derivation time was \JsonifiableDerivationMeanSeconds{} seconds, the median was
\JsonifiableDerivationMedianSeconds{} seconds, the 90th percentile was
\JsonifiableDerivationPNineZeroSeconds{} seconds, and the maximum was
\JsonifiableDerivationMaxSeconds{} seconds.

\begin{table}
  \centering
  \caption{Slowest successful derivations in the broad benchmark.}
  \begin{tabular}{llr}
    \toprule
    Declaration & Kind & Seconds \\
    \midrule
    \JsonifiableSlowestRows
    \bottomrule
  \end{tabular}
\end{table}

\subsection{Extracted Code Performance}

The extraction benchmark compares generated definitions against handwritten
Rocq definitions after both have been extracted to OCaml.  This keeps the
comparison in the same source language and extraction pipeline.  The benchmark
covers several shapes: a 256-constructor enumeration, a tuple-like constructor,
a mixed sum with both nullary and field-bearing constructors, and a recursive
binary tree.  Each row reports \texttt{JSONIFIABLE\_EXTRACTION\_RUNS} samples,
which defaults to 10.  The enum benchmark serializes and deserializes every
constructor; the other benchmarks run many iterations over representative
values.  These benchmarks are still small, but they exercise different
generated-code paths and are more informative than a single stress case.

\begin{table}
  \centering
  \caption{Extracted round-trip benchmarks.}
  \begin{tabular}{lrrrrr}
    \toprule
    Benchmark & Samples & Mean & Median & Min & Max \\
    \midrule
    \JsonifiableExtractionRows
    \bottomrule
  \end{tabular}
\end{table}

The intended conclusion is not that these microbenchmarks predict every
serializer workload.  Rather, they check that the deriver does not impose an
obvious extraction penalty on common shapes and on a case that stresses
constructor dispatch.

\section{Artifact}

The artifact is part of the contribution.  Tool papers benefit from evaluation
that reviewers can rerun without reconstructing informal commands from prose.
Our artifact provides a Docker path and a native path.  The main script rebuilds
the project from a clean state, runs the broad applicability benchmark, runs the
extraction benchmark, and regenerates the LaTeX macros included in this paper.

The generated files are meant to be read directly.  The JSON summary gives the
aggregate result, the CSV files contain row-level data, and the Markdown failure
report explains each failure category with representative Rocq diagnostics.
This setup makes the boundary of the tool explicit: reviewers can see both what
works and what is outside the supported fragment.

\section{Limitations}

The current deriver targets computational data in \texttt{Type} or
\texttt{Set}.  It does not attempt to serialize arbitrary propositions,
dependent families, coinductive values, function fields, or records whose field
types depend on earlier fields.  Some nested recursion through standard
containers is supported, but general nested and mutually recursive types remain
outside the current design.  This is a deliberate usability choice: partial
support for a few accidental nested forms can be worse than a clear diagnostic
when users cannot predict which structurally similar declarations will work.

The JSON representation is also intentionally modest.  The paper evaluates the
derivation method and the certified round-trip property, not complete adherence
to the JSON RFC or interoperability with every external JSON parser.

\section{Conclusion}

\texttt{rocq-json}'s ELPI deriver demonstrates that useful certified
serialization support can be built with a small user-facing interface and an
honest artifact.  The generated instances combine executable encoders and
decoders with Rocq-checked invertibility proofs, apply across a meaningful
fragment of existing library declarations, and extract to competitive code on a
large-enumeration stress test.  The engineering is specific to JSON, but the
pattern is broader: use ELPI to inspect declarations, generate computational
content and proof obligations together, and evaluate both applicability and
runtime behavior as first-class artifact outputs.

\end{document}
